\documentclass{article}

% Requires neurips_2026.sty from the ATTRIB CFP template link, in this directory.
\PassOptionsToPackage{square,numbers}{natbib}
\usepackage[dblblindworkshop,final]{neurips_2026}
\workshoptitle{Attributing Model Behavior at Scale: Data Attribution and Provenance}

\usepackage[utf8]{inputenc}
\usepackage[T1]{fontenc}
\usepackage{hyperref}
\usepackage{url}
\usepackage{booktabs}
\usepackage{amsmath}
\usepackage{placeins}

\title{Ablate-and-Measure Breaks When the Measurement Depends on the Ablation}

\author{%
  Abhimanyu Prasad \\
  \texttt{abhiprd2000@gmail.com}
}

\begin{document}

\maketitle

\begin{abstract}
Attribution methods are validated by ablation: remove a component, re-measure model
behavior, and check whether the predicted change occurred. This loop assumes the
measurement is available whether or not the ablation was applied. On agentic
benchmarks that assumption fails. Per-episode metrics are defined only for episodes
that reach a scoreable state, and ablations change how often episodes get there. We
document this with a 144-episode context-ablation study on $\tau^2$-bench that we could
not interpret. The fraction of episodes yielding a defined grounding rate ranges from
50.0\% to 91.7\% across conditions; for the benchmark's own action-level score it
ranges from 4.2\% to 33.3\%. Because paired testing retains only cells defined in both
arms, our \emph{control} condition's measured rate ranges from 0.096 to 0.304 depending
on which treatment it is compared against, a 21-point swing in a condition that never
changed, exceeding the effect the study was powered to detect. One treatment reverses
sign between two defensible analyses of the same episodes. We report this as a failed
experiment with a transferable lesson: as data attribution extends to agentic
behaviors, the metrics being predicted are exactly the outcome-gated kind that breaks
this loop, and four cheap diagnostics reveal whether it is happening.
\end{abstract}

\section{The claim}

Contributive attribution asks which data caused a model behavior, and the standard way
to check an attribution method is counterfactual: remove the implicated data, retrain
or ablate, measure the behavior, and see whether the method predicted the change.
Leave-one-out correlation \citep{koh2017influence}, datamodels
\citep{ilyas2022datamodels}, and the linear datamodeling score \citep{park2023trak} all
score a method by how well it predicts a model output function $f$ under perturbations
of the training set.

The construction has a silent precondition: it presumes $f$ is defined on every
perturbed run. For a next-token probability or an accuracy on a fixed test set this is
unremarkable, since the metric exists for every input by construction.

It stops being unremarkable when the measured behavior is agentic. A tool-using agent
produces a score only if its episode terminates in a state the evaluator will score.
Episodes that exhaust a step budget, accumulate tool errors, or overflow a context
window yield no measurement at all. They are undefined, not zero. And whether an
episode reaches a scoreable state is itself affected by the ablation.

When that happens the two arms are no longer the same population. The analysis
conditions on an outcome downstream of the treatment, which is conditioning on a
collider \citep{elwert2014endogenous}. Our contribution is not to name the phenomenon
but to show, with measurements, that it is severe enough in agentic evaluation to
invert a conclusion, and to propose what to report so that readers can see it.

We arrived here by failing. We set out to attribute tool-call failures to parts of an
agent's context, ran the full grid, and could not interpret the result. We think the
failure is more useful than the experiment would have been.

\section{What we ran}

$\tau^2$-bench \citep{barres2025tau2}, retail domain, pinned at the v1.0.1 release
commit, which incorporates the benchmark's task-quality corrections. Our attrition is
therefore not attributable to the task-specification errors those corrections
addressed. Agent and simulated user are both Qwen2.5-7B-Instruct-AWQ under vLLM at temperature 0 with a
200-step cap. Using one model on both sides fixes the user simulator so that
differences are attributable to the agent's context. Ablation is an established
practice on this benchmark; the original paper reports several.

Six conditions: a full-context control and five ablations removing, respectively, the
entire policy document; all policy sections except those governing the task's actions;
all tool descriptions; conversation history beyond the last three messages; and
policy-framing sentences in the system instruction. Each was implemented by wrapping
the episode runner and verified by dumping and diffing the served chat-completion
payload for every task. Each condition changed only what it claimed to
(Appendix~\ref{app:verify}).

The grid was 12 tasks $\times$ 2 seeds $\times$ 6 conditions, or 144 cells, of which
138 episodes completed. The primary metric was \emph{grounding rate}: the fraction of
identifier-valued tool-call arguments appearing in prior context rather than being
fabricated. Fabricated identifiers are a documented failure mode on this benchmark
family, and the original $\tau$-bench paper reports per-task rates of tool calls
carrying non-existent user, product, order, or item IDs \citep{yao2024taubench}.
Analysis was pre-registered before the grid ran: paired sign-flip permutation test,
$\alpha = 0.05$, Holm correction across five conditions, and a 20-point falsification
threshold matched to the grid's minimum detectable effect at 87\% power.

\section{Definedness is condition-dependent}

We chose grounding rate precisely to avoid the benchmark's own scoring gate. That gate
is explicit in source: the evaluator returns a null result without reaching action
checking whenever an episode does not terminate in one of two clean stop states. Its
action-level score consequently exists for between 1 and 8 of 24 episodes per
condition, or 4.2\% to 33.3\% (Appendix~\ref{app:acc}). One condition retains a single
defined cell and therefore drops out of every paired contrast without appearing to.
A further 21 episodes lost that score to a fault in the assertion judge on our side,
again unevenly across conditions, from 1 to 6. That fault was ours, not the benchmark's,
and touches only this secondary metric: the grounding rate is computed from transcripts
and never invokes the evaluator. A measurement can go missing for reasons unrelated to
agent behavior, and an aggregated score does not reveal it.

Grounding rate avoids that gate but not the mechanism behind it. It is undefined when
an episode makes no identifier-bearing tool calls, or ends before its grounding pool
fills, and that too depends on condition:

\begin{table}[h]
\centering
\caption{Episodes yielding a defined grounding rate, by condition.}
\begin{tabular}{lrr}
\toprule
Condition & Defined (of 24) & Share \\
\midrule
Control (Full Context) & 22 & 91.7\% \\
Policy Truncated & 22 & 91.7\% \\
History Truncated & 21 & 87.5\% \\
Schemas Stripped & 20 & 83.3\% \\
System Minimized & 13 & 54.2\% \\
Policy Removed & 12 & 50.0\% \\
\bottomrule
\end{tabular}
\end{table}

Half the policy-removed episodes produce no primary measurement. Choosing a metric
outside the benchmark's scoring path reduced the problem without removing it: both
metrics require the episode to get far enough to produce the thing being measured.

The driver is visible in termination reasons, where clean terminations range from 4 of
24 to 11 of 24 across conditions, a 29.1-point spread (Appendix~\ref{app:term}). Our
pre-registered rule declared the primary confounded above a 10-point spread. It fired.

Reward-design defects in this benchmark family are documented rather than hypothetical
\citep{zhu2025abc}, as were task-specification errors in its predecessor's airline
domain \citep{cuadron2026saber}, which prompted the corrections in the release we use.
Our point is different and narrower: even a correctly
specified reward creates a selection problem if whether it is computed depends on the
manipulation. Nor is the practice hypothetical: recent work analyzing false success on
this same benchmark discards trajectories with missing rewards as routine
preprocessing \citep{advani2026silent}, a step that is standard, reasonable, and
exactly the one whose innocuousness we question.

We note that the error-terminated episodes are behavioral, not harness defects. The
agent calls real tools with well-formed arguments containing fabricated identifiers
until the error budget is exhausted. We verified that tool-call parsing was configured
correctly and that every failing call was structurally valid.

\section{The control moves}

Definedness rates alone do not convey how much this matters. This does.

Our test pairs each treatment cell with the control cell at the same (task, seed),
which is standard and reduces task variance. But pairing retains only cells defined in
\emph{both} arms, so each treatment induces a different retained subset of the control. The control never changed. Its measured value does:

\begin{table}[h]
\centering
\caption{The control condition's own grounded rate, computed on each treatment's
retained pairs. Its marginal rate over all 22 defined cells is 0.289.}
\begin{tabular}{lrr}
\toprule
Paired Against & Pairs & Control's Grounded Rate \\
\midrule
Policy Removed & 11 & 0.096 \\
Schemas Stripped & 18 & 0.198 \\
Policy Truncated & 20 & 0.265 \\
System Minimized & 11 & 0.299 \\
History Truncated & 19 & 0.304 \\
\bottomrule
\end{tabular}
\end{table}

Across pairings the control spans 21 points, larger than the 20-point effect the grid
was designed to detect, and produced entirely by which cells survived.

The consequence is a sign reversal. For the policy-removed condition, which cuts the
served system prompt from 7{,}056 characters to 357:

\begin{itemize}
\item \textbf{Marginal:} ungrounded rate 71.7\% against the control's 71.1\%. Removing
the policy document makes grounding \emph{0.6 points worse}.
\item \textbf{Paired:} grounded rate 0.309 against the control's 0.096 on the 11 shared
cells. Removing it makes grounding \emph{21.2 points better} (raw $p = 0.047$,
Holm-adjusted $p = 0.235$).
\end{itemize}

Same episodes, same metric, opposite conclusion. Neither is a coding error. They
describe different populations, and nothing in either table tells a reader which
population they are looking at.

For completeness: under the pre-registered test no condition is significant after
correction, with a minimum adjusted $p$ of 0.235, and none crosses the 20-point
threshold under worst-case or best-case imputation of lost episodes
(Appendix~\ref{app:primary}). At 87\% power for 20 points we can exclude large effects
and nothing smaller, and even that bounded null describes the observed subpopulation
rather than the intended one.

\section{Why this generalizes to data attribution}

We ablated context, not training data, and do not claim otherwise. The transferable
part is the validation loop, which is shared.

Evaluating a contributive attribution method means predicting how a behavior changes
when implicated data is removed, then removing it and checking. Where the behavioral
measurement is defined for every input, such as a loss, a probability, or an accuracy
on a fixed set, the loop is sound. As attribution moves toward agentic behaviors, which
several of this workshop's motivating applications require, such as auditing what data
caused an agent to take an action or tracing a tool-use failure to its provenance, the
measured quantity becomes an episode-level outcome of exactly the gated kind we
document.

At that point the counterfactual check inherits the problem. Removing training data may
change how often an agent's episodes terminate scoreably. The retained comparison set
then differs between arms, the attribution method's predictions are scored against a
shifting population, and the method can be credited or blamed for a selection effect.
Nothing about the attribution method is wrong; the harness that validates it measures
something other than what it reports.

\section{Four things to report}

The problem is cheap to detect. We propose that any ablation-based evaluation of
agentic behavior report, alongside its headline table:

\begin{enumerate}
\item \textbf{Per-condition metric definedness.} How many episodes produced a value,
not just the mean of those that did. A condition contributing half as many measurements
as the control is itself a result.
\item \textbf{Per-condition termination distribution.} Clean stops, step exhaustion,
error exhaustion, context overflow, and crashes. This is where the selection
originates.
\item \textbf{Control stability under pairing.} If the analysis is paired, report the
control's value on each retained subset. If it moves, the comparison is not what it
appears to be.
\item \textbf{Bounds rather than exclusions.} Report the primary with dropped episodes
imputed at both extremes. If the bounds straddle the decision threshold, say so.
\end{enumerate}

These complement rather than duplicate existing guidance. The Agentic Benchmark
Checklist \citep{zhu2025abc} asks whether a reward is \emph{correct}; our diagnostics
ask whether it is \emph{defined}, and whether definedness depends on the manipulation
under test. None require extra compute: all four derive from records a benchmark run
already produces, and we report all four for our own grid.

\section{Limitations}

One 7B model, one benchmark domain, 12 tasks and 2 seeds. The mechanism is structural
and we expect it to recur, but its magnitude here is specific to this setup.

The confound applies to our own null. Section 4's result inherits the problem described
in Section 3. We report it as evidence about the design, not about the hypothesis.

Two deviations from our pre-registration, disclosed rather than smoothed over. The rule
for scoring error-terminated episodes was written after pilot-level accuracy had been
seen; it was derived by symmetry with an existing frozen rule, but it was not written
blind. The per-cell aggregation function required by the frozen primary test did not
exist when the plan was frozen, and was implemented after marginal point estimates had
been seen; the metric, test, $\alpha$, and correction family were unchanged. The
combined attrition statistic that triggered our confound rule was operationalized after
the fact, so we publish raw per-condition termination counts for readers who prefer a
different definition.

Our truncated-policy condition selects sections using ground-truth actions. It
upper-bounds what perfect retrieval would provide and does not simulate a retriever.


\begin{thebibliography}{9}

\bibitem{advani2026silent}
L.~Advani.
\newblock From confident closing to silent failure: Characterizing false success in
LLM agents.
\newblock \emph{arXiv:2606.09863}, 2026.

\bibitem{barres2025tau2}
V.~Barres, H.~Dong, S.~Ray, X.~Si, and K.~Narasimhan.
\newblock $\tau^2$-bench: Evaluating conversational agents in a dual-control
environment.
\newblock \emph{arXiv:2506.07982}, 2025.

\bibitem{cuadron2026saber}
A.~Cuadron, P.~Yu, Y.~Liu, and A.~Gupta.
\newblock SABER: Small actions, big errors. Safeguarding mutating steps in LLM agents.
\newblock In \emph{ICLR 2026 Workshop on Memory for LLM-Based Agentic Systems}, 2026.

\bibitem{elwert2014endogenous}
F.~Elwert and C.~Winship.
\newblock Endogenous selection bias: The problem of conditioning on a collider
variable.
\newblock \emph{Annual Review of Sociology}, 40:31--53, 2014.

\bibitem{ilyas2022datamodels}
A.~Ilyas, S.~M.~Park, L.~Engstrom, G.~Leclerc, and A.~Madry.
\newblock Datamodels: Predicting predictions from training data.
\newblock \emph{arXiv:2202.00622}, 2022.

\bibitem{koh2017influence}
P.~W.~Koh and P.~Liang.
\newblock Understanding black-box predictions via influence functions.
\newblock In \emph{ICML}, 2017.

\bibitem{park2023trak}
S.~M.~Park, K.~Georgiev, A.~Ilyas, G.~Leclerc, and A.~Madry.
\newblock TRAK: Attributing model behavior at scale.
\newblock In \emph{ICML}, volume 202 of \emph{PMLR}, pages 27074--27113, 2023.

\bibitem{yao2024taubench}
S.~Yao, N.~Shinn, P.~Razavi, and K.~Narasimhan.
\newblock $\tau$-bench: A benchmark for tool-agent-user interaction in real-world
domains.
\newblock \emph{arXiv:2406.12045}, 2024.

\bibitem{zhu2025abc}
Y.~Zhu, T.~Jin, Y.~Pruksachatkun, A.~Zhang, S.~Liu, S.~Cui, S.~Kapoor, S.~Longpre,
K.~Meng, R.~Weiss, F.~Barez, R.~Gupta, J.~Dhamala, J.~Merizian, M.~Giulianelli,
H.~Coppock, C.~Ududec, J.~Sekhon, J.~Steinhardt, A.~Kellermann, S.~Schwettmann,
M.~Zaharia, I.~Stoica, P.~Liang, and D.~Kang.
\newblock Establishing best practices for building rigorous agentic benchmarks.
\newblock \emph{arXiv:2507.02825}, 2025.

\end{thebibliography}

\appendix

\section{Ablation verification}
\label{app:verify}

\begin{table}[!ht]
\centering
\small
\caption{Served payload per condition, verified on all tasks by diffing against the
control. Policy-truncated length varies by task.}
\begin{tabular}{lrrrc}
\toprule
Condition & System Prompt (chars) & Tools & Tools w/ Descriptions & History Sent \\
\midrule
Full Context & 7056 & 16 & 16 & Full \\
Policy Removed & 357 & 16 & 16 & Full \\
Policy Truncated & 2311 to 4621 & 16 & 16 & Full \\
Schemas Stripped & 7056 & 16 & 0 & Full \\
History Truncated & 7056 & 16 & 16 & Last 3 \\
System Minimized & 6913 & 16 & 16 & Full \\
\bottomrule
\end{tabular}
\end{table}

\FloatBarrier
\section{Termination distribution}
\label{app:term}

\begin{table}[!ht]
\centering
\small
\caption{Termination reason counts over all 144 cells. Crashes account for the 6
episodes that did not complete.}
\begin{tabular}{lrrrrr}
\toprule
Condition & Normal & Agent Error & Step Cap & Context Overflow & Crash \\
\midrule
Full Context & 11 & 9 & 3 & 0 & 1 \\
Policy Removed & 11 & 7 & 4 & 0 & 2 \\
Policy Truncated & 11 & 4 & 8 & 0 & 1 \\
Schemas Stripped & 10 & 11 & 3 & 0 & 0 \\
History Truncated & 5 & 17 & 2 & 0 & 0 \\
System Minimized & 4 & 9 & 6 & 3 & 2 \\
\bottomrule
\end{tabular}
\end{table}

\FloatBarrier
\section{Pre-registered primary test}
\label{app:primary}

Positive difference means the condition grounded worse than control.

\begin{table}[!ht]
\centering
\small
\caption{Paired sign-flip permutation test on grounded rate, Holm-corrected across five
conditions.}
\begin{tabular}{lrrrrrr}
\toprule
Condition & Pairs & Control & Treatment & Diff & Raw $p$ & Holm $p$ \\
\midrule
Policy Removed & 11 & 0.096 & 0.309 & $-0.212$ & 0.047 & 0.235 \\
Policy Truncated & 20 & 0.265 & 0.240 & $+0.025$ & 0.816 & 1.000 \\
Schemas Stripped & 18 & 0.198 & 0.261 & $-0.063$ & 0.485 & 1.000 \\
History Truncated & 19 & 0.304 & 0.180 & $+0.123$ & 0.135 & 0.542 \\
System Minimized & 11 & 0.299 & 0.208 & $+0.090$ & 0.517 & 1.000 \\
\bottomrule
\end{tabular}
\end{table}

Crash-exclusion bounds, given as change against control under worst-case and best-case
imputation: policy removed $+3.4$ and $-6.6$; policy truncated $+1.4$ and $+1.4$;
schemas stripped $+0.8$ and $+5.1$; history truncated $+7.8$ and $+12.2$; system
minimized $+10.7$ and $-12.7$. No condition's bounds cross the 20-point threshold.

\section{Action-level score definedness}
\label{app:acc}

\begin{table}[!ht]
\centering
\small
\caption{Episodes yielding a defined action-level score, out of 24 per condition.
``Not Reached'' covers episodes whose termination reason prevented action checking.
``Evaluation Failed'' covers episodes lost to a fault in the assertion judge on our
side. Counts derive from the evaluator note stored in each released record.}
\begin{tabular}{lrrr}
\toprule
Condition & Scored & Not Reached & Evaluation Failed \\
\midrule
Policy Removed & 8 & 13 & 3 \\
Policy Truncated & 8 & 13 & 3 \\
Full Context & 7 & 13 & 4 \\
Schemas Stripped & 4 & 14 & 6 \\
System Minimized & 3 & 20 & 1 \\
History Truncated & 1 & 19 & 4 \\
\bottomrule
\end{tabular}
\end{table}

\section{Marginal grounding rates}

Ungrounded argument rate, macro-averaged over episodes with a defined rate: full
context 71.1\% ($n{=}22$), policy removed 71.7\% (12), policy truncated 72.6\% (22),
schemas stripped 73.2\% (20), system minimized 76.6\% (13), history truncated 80.2\%
(21). The agent fabricates roughly three of every four identifier arguments in every
condition, including the control.

\section{Power}

Minimum detectable effect on grounding rate at 12 tasks and 2 seeds, by simulation with
60 replicates per point: 5 points gives 0\% power, 10 points 8\%, 15 points 43\%, 20
points 87\%, and 30 points 100\%. The 20-point threshold was set to match this curve
before the grid ran.

\section{Reproducibility}

Benchmark pinned at the v1.0.1 release commit. Both models are Qwen2.5-7B-Instruct-AWQ under vLLM
at temperature 0, with a 200-step cap, on two T4 GPUs, for roughly 16 GPU-hours across
two sessions. All 144 episode records, the payload verification dump, and the analysis
plan with its frozen and archived sections are available at
\url{https://github.com/abhiprd2000/attrib-context-ablation}.

\end{document}